\documentclass[12pt,a4paper]{article}
\usepackage[utf8]{inputenc}
\usepackage[spanish]{babel}
\usepackage{graphicx}
\usepackage{amsmath}
\usepackage{amssymb}
\usepackage{booktabs}
\usepackage{hyperref}

\title{Capítulo: Integración Neuromórfica en Hardware Memristivo para Control Robótico Autónomo mediante R-STDP}
\author{Yamil Uchani}
\date{\today}

\begin{document}

\maketitle

\begin{abstract}
El diseño de algoritmos de Aprendizaje por Refuerzo (RL) adaptados a la robótica de borde exige reducir drásticamente los costos energéticos y computacionales de la retropropagación convencional. En este capítulo se presenta el desarrollo, modelado físico y validación del simulador neuromórfico \textit{NeuroBot}, un sistema de control adaptativo basado en una red memristiva \textit{single-layer} de $24 \times 12$ acoplada a un entorno robótico dinámico (CartPole). Mediante una regla de aprendizaje por picos modulada por recompensa (R-STDP), resolución de circuitos con Análisis Nodal Modificado (MNA) y una exploración dinámica por temperatura adaptativa, el agente logró convergencia a la recompensa máxima promedio ($\bar{R} \approx 487.7$) en menos de 30 episodios. Se presenta un estudio de ablación exhaustivo ($A/B/C/D$) y pruebas de reproducibilidad en múltiples semillas para fundamentar la física del dispositivo y la asignación de crédito temporal.
\end{abstract}

\section{Introducción y Motivación}
El paradigma de cómputo in-memory neuromórfico permite realizar multiplicaciones vector-matriz de forma analítica en leyes físicas (Leyes de Ohm y Kirchhoff), eliminando el cuello de botella de Von Neumann. Sin embargo, aplicar redes memristivas al control de sistemas dinámicos continuos requiere superar tres desafíos metodológicos:
\begin{enumerate}
    \item \textbf{Asignación de crédito temporal:} Conectar recompensas ambientales escalares retrasadas con cambios locales de conductancia a nivel de celda sin retropropagación.
    \item \textbf{No idealidades físicas del dispositivo:} Considerar los efectos de caída de línea ($R_{\text{wire}}$), corrientes parásitas (Sneak Paths) y no linealidades de programación memristiva (modelo de Strukov/Prezioso).
    \item \textbf{Trampas locales en el espacio de estados:} Evitar la saturación prematura de conductancia mediante la modulación adaptativa de exploración/explotación.
\end{enumerate}

\section{Arquitectura del Sistema ($24 \times 12$)}
El cerebro neuromórfico del agente consiste en un crossbar memristivo de $N=24$ filas de entrada y $M=12$ columnas de salida.

\subsection{Codificador de Estados (Encoder Top-K)}
El vector de observación continuo $\mathbf{s} = [x, \dot{x}, \theta, \dot{\theta}] \in \mathbb{R}^4$ se proyecta mediante una batería de receptores gaussianos distribuidos a lo largo del dominio de cada variable. Se aplica un mecanismo Top-K ($K=10$) que selecciona las 10 filas de mayor activación, generando un patrón binario disperso de voltajes de lectura $V_{\text{rows}} \in \{0, V_{\text{read}}\}$.

\subsection{Crossbar Memristivo y Solución Nodal (MNA)}
La corriente integrada por columna $\mathbf{I}_{\text{col}} \in \mathbb{R}^{12}$ se calcula resolviendo la matriz MNA con resistencia de línea $R_{\text{wire}} = 0.5\,\Omega$ y resistencia de sensado $R_{\text{sense}} = 100\,\Omega$:
\begin{equation}
    \mathbf{I}_{\text{col}} = \text{MNA}(\mathbf{G}, V_{\text{rows}}, R_{\text{wire}}, R_{\text{sense}})
\end{equation}
donde $G_{ij} = \frac{1}{R_{\text{on}} + X_{ij}(R_{\text{off}} - R_{\text{on}})}$ representa la conductancia de la celda en el estado interno $X_{ij} \in [0, 1]$.

\subsection{Decodificador y Acción (Softmax e Inercia de Motor)}
Las 12 columnas se agrupan en pares de conductancias asociadas a 6 niveles discretos de fuerza activa para giro hacia la izquierda ($1 - 30\,\text{N}$) y 6 para la derecha ($1 - 30\,\text{N}$). La probabilidad de selección de acción $\pi_j$ sigue una distribución Softmax parametrizada por la temperatura ambiental $T$:
\begin{equation}
    \pi_j = \frac{\exp(I_j / T)}{\sum_{k=1}^M \exp(I_k / T)}
\end{equation}
Se aplica una tasa de cambio (slew rate) de $50/50$ en la dinámica del actuador para simular la inercia del motor robótico real.

\section{Regla de Aprendizaje R-STDP y Traza Volátil}
El aprendizaje ocurre al final de cada episodio según la regla de gradiente de política modulada por una línea base móvil $b$:
\begin{equation}
    \Delta X_{ij} = -\eta \cdot E_{ij} \cdot (R_{\text{ep}} - b)
\end{equation}
donde $E_{ij} = \sum_{t} \text{pre}_i(t) \cdot (\mathbf{1}_{j = a(t)} - \pi_j(t))$ es la traza de elegibilidad acumulada durante el episodio. Físicamente, esta traza es replicada en tiempo real mediante un arreglo gemelo de memristores volátiles (dispositivos de relajación temporal).

Para preservar la estabilidad de la matriz, la actualización $\Delta X$ utiliza recorte por norma de gradiente ($\max|\Delta X| \le 0.05$) y salta la reescritura si $R_{\text{ep}} \ge 430$, evitando degradación sobre políticas óptimas alcanzadas.

\section{Resultados Experimentales y Discusión}

\subsection{Convergencia y Reproducibilidad Multi-Semilla}
La Figura~\ref{fig:curva} ilustra la curva de aprendizaje en 150 episodios para 3 semillas independientes (Seed 0, Seed 1 y Seed 2).

\begin{figure}[htbp]
    \centering
    \includegraphics[width=0.85\textwidth]{figuras/curva_aprendizaje_24x12.pdf}
    \caption{Curva de aprendizaje del agente en arquitectura $24 \times 12$. Muestra convergencia ultrarrápida a $\bar{R} \approx 487.7$ en $<30$ episodios y la banda de confianza $\pm 1\sigma$ sobre múltiples semillas.}
    \label{fig:curva}
\end{figure}

El sistema converge a $\bar{R} \approx 487$ en 3 de 5 semillas en $<30$ episodios, representando una aceleración de $35\times$ en comparación con la arquitectura previa de $12 \times 4$. Las semillas restantes caen en mínimos locales ($\bar{R} \approx -20$), lo cual es atribuible a la variabilidad de la proyección aleatoria inicial del encoder. Esta tasa de éxito del $60\%$ refleja el comportamiento característico del algoritmo R-STDP en ausencia de un buffer de experiencia.

\subsection{Estudio de Ablaciones}
Para validar la contribución de cada bloque del agente, se evaluaron cuatro configuraciones (Figura~\ref{fig:ablaciones}):
\begin{itemize}
    \item \textbf{A (Control Base):} Sistema completo con decaimiento térmico y traza acumulativa ($\bar{R} \approx 491.9$ en semillas exitosas, $60\%$ de tasa de convergencia en 5 semillas aleatorias).
    \item \textbf{B (Sin Adaptación Térmica, $T=0.10$ Fijo):} Colapso por congelamiento estocástico ($\bar{R} = -39.0$). Demuestra que sin alta exploración inicial, el agente congela la búsqueda de políticas.
    \item \textbf{C (Sin Traza de Elegibilidad, $\varepsilon=0$):} Colapso total ($\bar{R} = -33{,}3$, $\Delta X=0$), evidenciando que la asignaci\'{o}n de cr\'{e}dito temporal es estrictamente indispensable para R-STDP.
    \item \textbf{D (Hebbiano Sin Modulación $R_{\text{ep}}-b$):} Divergencia y colapso ($\bar{R} = -39.0$) por falta del escalar de refuerzo global.
\end{itemize}

\begin{figure}[htbp]
    \centering
    \includegraphics[width=0.85\textwidth]{figuras/comparativa_ablaciones_24x12.pdf}
    \caption{Estudio de ablación comparativo (A, B, C, D) en el crossbar $24 \times 12$.}
    \label{fig:ablaciones}
\end{figure}

\subsection{Heatmap de la Matriz Memristiva Consolidada ($X$)}
La Figura~\ref{fig:heatmap} muestra la distribución final del estado memristivo $X \in [0, 1]$. Se observa una estructura bimodal no saturada ($\mu = 0.165, \sigma = 0.118$) con bandas claras de activación sináptica diferenciadas por la velocidad y el ángulo del poste.

\begin{figure}[htbp]
    \centering
    \includegraphics[width=0.65\textwidth]{figuras/heatmap_matriz_X_24x12.pdf}
    \caption{Heatmap del estado interno de conductancias $X$ tras la convergencia del agente.}
    \label{fig:heatmap}
\end{figure}


\section{Validación del Modelo Físico del Crossbar}

Para certificar la fidelidad del simulador neuromórfico respecto al hardware real (configurando $R_{\text{sense}} = 10^{-3}\,\Omega$ en el crossbar aislado), se ejecutó una suite de 10 experimentos autónomos. Cada test es independiente, reproducible y genera sus propias métricas y figuras.

\subsection{Exp.~1 --- Corrientes parásitas vs.\ tamaño (sneak ratio)}
La Figura~\ref{fig:exp01} muestra el crecimiento monótono del \textit{sneak ratio} al escalar el arreglo de $N=4$ ($0.34\%$) a $N=32$ filas ($2.44\%$, con $M=12$ fijo).
\begin{figure}[htbp]
  \centering
  \includegraphics[width=0.65\textwidth]{experiments/figures/exp01_sneak_vs_N.png}
  \caption{Sneak ratio vs.\ tamaño. Crecimiento monótono confirma el efecto parásito escalar.}
  \label{fig:exp01}
\end{figure}

\subsection{Exp.~2 --- MNA real vs.\ modelo ideal}
La Figura~\ref{fig:exp02} compara corrientes MNA e ideales para cuatro tamaños. A $N=24$ el error relativo medio alcanza el $3.63\%$, demostrando la necesidad de resolver las ecuaciones circuitales completas.
\begin{figure}[htbp]
  \centering
  \includegraphics[width=0.90\textwidth]{experiments/figures/exp02_mna_vs_ideal.png}
  \caption{Diagrama de dispersión MNA vs.\ Ideal. El error crece de forma no lineal con $N$.}
  \label{fig:exp02}
\end{figure}

\subsection{Exp.~3 --- Caída de tensión en línea ($R_{\text{wire}}$)}
La Figura~\ref{fig:exp03} ilustra cómo $R_{\text{wire}}$ distorsiona el perfil de corriente. Para $R_{\text{wire}} \ge 1\,\Omega$ las variaciones entre columnas llegan a $\pm 30\%$.
\begin{figure}[htbp]
  \centering
  \includegraphics[width=0.72\textwidth]{experiments/figures/exp03_rwire_effect.png}
  \caption{Perfil de $I_{\text{col}}$ para distintos $R_{\text{wire}}$. Mayor resistencia implica mayor no-uniformidad.}
  \label{fig:exp03}
\end{figure}

\subsection{Exp.~4 --- Precisión en VMM}
La Figura~\ref{fig:exp04} muestra la distribución del error relativo en VMM para 1000 vectores aleatorios. El error medio crece de $0.37\%$ ($N=4$) a $3.14\%$ ($N=24$).
\begin{figure}[htbp]
  \centering
  \includegraphics[width=0.70\textwidth]{experiments/figures/exp04_vmm_precision.png}
  \caption{Histograma del error relativo en VMM por tamaño de crossbar.}
  \label{fig:exp04}
\end{figure}

\subsection{Exp.~5 --- Half-select disturb: V/2 vs.\ V/3}
La Figura~\ref{fig:exp05} compara la perturbación física en celdas half-selected. V/2 produce un disturb del $50.0\%$ (selectividad $2.0\times$); V/3 lo reduce al $33.3\%$ (selectividad $3.0\times$), en exacta concordancia con el modelo lineal de Strukov.
\begin{figure}[htbp]
  \centering
  \includegraphics[width=0.95\textwidth]{experiments/figures/exp05_v2_vs_v3.png}
  \caption{Evolución de $X$ para celda fully-, half- y no-selected bajo 3 pulsos.}
  \label{fig:exp05}
\end{figure}

\subsection{Exp.~6 --- Retención de estado (no volatilidad)}
Tras 1000 pasos de lectura con $V_{\text{rows}}=0\,\text{V}$, el cambio máximo en $X$ es $\Delta_{\max} = 0$, certificando retención perfecta del dispositivo no volátil.
\begin{figure}[htbp]
  \centering
  \includegraphics[width=0.62\textwidth]{experiments/figures/exp06_retention.png}
  \caption{Media de $X$ durante 1000 pasos en reposo. Variación nula = retención perfecta.}
  \label{fig:exp06}
\end{figure}

\subsection{Exp.~7 --- Linealidad LTP/LTD}
La Figura~\ref{fig:exp07} muestra curvas monótonas y acotadas en $[0,1]$ bajo pulsos de potenciación (LTP, $+1\,\text{V}$) y depresión (LTD, $-1\,\text{V}$). La celda responde progresivamente alcanzando $x=1.000$ y $x=0.000$.
\begin{figure}[htbp]
  \centering
  \includegraphics[width=0.88\textwidth]{experiments/figures/exp07_ltp_ltd.png}
  \caption{LTP/LTD y derivada discreta. La no-linealidad satura suavemente los extremos.}
  \label{fig:exp07}
\end{figure}

\subsection{Exp.~8 --- Variabilidad dispositivo a dispositivo (D2D)}
La Figura~\ref{fig:exp08} cuantifica la dispersión de $I_{\text{col}}$ entre 20 crossbars virtuales. El CV normalizado por el rango dinámico de conductancia es $141.3\%$, reflejando una dispersión inter-dispositivo clara y físicamente realista.
\begin{figure}[htbp]
  \centering
  \includegraphics[width=0.88\textwidth]{experiments/figures/exp08_d2d_reproducibility.png}
  \caption{Dispersión de $I_{\text{col}}$ y variación de $\bar{G}_{\text{mean}}$ entre 20 crossbars virtuales.}
  \label{fig:exp08}
\end{figure}

\subsection{Exp.~9 --- Costo computacional del solver MNA}
La Figura~\ref{fig:exp09} revela una complejidad empírica de $\mathcal{O}(N^{0.88})$, inferior al $\mathcal{O}(N^3)$ teórico. A $N=24$, el costo es de $1.96\,\text{ms/paso}$, compatible con simulación en tiempo real.
\begin{figure}[htbp]
  \centering
  \includegraphics[width=0.62\textwidth]{experiments/figures/exp09_computational_cost.png}
  \caption{Costo del solver MNA en log-log. Ajuste empírico $\mathcal{O}(N^{0.88})$.}
  \label{fig:exp09}
\end{figure}

\subsection{Exp.~10 --- Crosstalk entre columnas}
La Figura~\ref{fig:exp10} expone el efecto de crosstalk: al excitar únicamente la fila 0, las otras columnas muestran una desviación media del $0.75\%$ por interacción parásita MNA.
\begin{figure}[htbp]
  \centering
  \includegraphics[width=0.88\textwidth]{experiments/figures/exp10_crosstalk.png}
  \caption{$I_{\text{col}}$ real vs.\ ideal al excitar únicamente la fila 0.}
  \label{fig:exp10}
\end{figure}

\section{Conclusiones}
La validación integral del pipeline ---desde el solver MNA hasta la política de control--- cierra la brecha entre simuladores de dispositivos físicos y frameworks de Machine Learning. Los 10 experimentos de la suite confirman:
\begin{enumerate}
    \item El crecimiento monótono del sneak ratio con el tamaño del arreglo ($0.34\%$ a $2.44\%$);
    \item Las selectividades teóricas exactas ($2.0\times$ V/2, $3.0\times$ V/3) en el modelo lineal de Strukov;
    \item La retención perfecta del dispositivo no volátil ($\Delta_{\max} = 0$);
    \item La variabilidad realista Device-to-Device ($CV_{\text{norm}} = 141.3\%$); y
    \item La convergencia del agente R-STDP en $<30$ episodios con $\bar{R} \approx 488$ sobre el crossbar físico $24 \times 12$.
\end{enumerate}

\end{document}
